\documentclass[UTF8]{ctexart}
\usepackage{amsmath, amssymb, geometry}
\usepackage{xcolor}
\usepackage{enumitem}
\usepackage{tcolorbox}
\usepackage{cases}
\usepackage{fancyhdr}
\geometry{a4paper, margin=1in}

% 定义红色环境
\newenvironment{redenv}{\color{red}}{}

% 定义详细解析环境
\newenvironment{detailedanalysis}{%
    \color{red}
    \begin{enumerate}[label=\textbf{第\arabic*步：}, leftmargin=*, align=left, itemsep=1.5ex]
}{%
    \end{enumerate}
}

% 定义概念框
\newenvironment{conceptbox}{%
    \begin{tcolorbox}[colback=red!5, colframe=red!50!black, title=概念解析, fonttitle=\bfseries]
}{%
    \end{tcolorbox}
}

% 定义公式推导环境
\newenvironment{formuladerivation}{%
    \begin{enumerate}[label={\textbf{公式\arabic*：}}, leftmargin=*, align=left]
}{%
    \end{enumerate}
}

% 定义题目和解析的分隔线
\newcommand{\problemrule}{\noindent\rule{\textwidth}{1.5pt}\vspace{-0.8em}}
\newcommand{\analysisrule}{\noindent\rule{\textwidth}{0.8pt}\vspace{-0.8em}}

% 宏定义：方便排版选择题选项
\newcommand{\choicesFour}[4]{
    \begin{enumerate}[label=\Alph*., itemsep=0pt, parsep=0pt, topsep=5pt, labelsep=0.5em]
        \begin{minipage}{0.22\linewidth} \item #1 \end{minipage}
        \begin{minipage}{0.22\linewidth} \item #2 \end{minipage}
        \begin{minipage}{0.22\linewidth} \item #3 \end{minipage}
        \begin{minipage}{0.22\linewidth} \item #4 \end{minipage}
    \end{enumerate}
}
\newcommand{\choicesTwo}[4]{
    \begin{enumerate}[label=\Alph*., itemsep=0pt, parsep=0pt, topsep=5pt, labelsep=0.5em]
        \begin{minipage}{0.45\linewidth} \item #1 \end{minipage}
        \begin{minipage}{0.45\linewidth} \item #2 \end{minipage}
        \begin{minipage}{0.45\linewidth} \item #3 \end{minipage}
        \begin{minipage}{0.45\linewidth} \item #4 \end{minipage}
    \end{enumerate}
}
\newcommand{\choices}[4]{
    \begin{enumerate}[label=\Alph*., itemsep=0pt, parsep=0pt, topsep=5pt, labelsep=0.5em]
        \item #1
        \item #2
        \item #3
        \item #4
    \end{enumerate}
}

\newcommand{\blank}[1]{\underline{\hspace{#1}}}

\begin{document}

\section*{第十二章 无穷级数提高十五题}

% ========== 第1题 ==========
\problemrule
\section*{【题目1】级数基本性质}
有下列关于数项级数的命题，其中正确的个数为 \blank{1cm}。
\begin{enumerate}[label=(\arabic*)]
    \item 若 $ \lim_{n\to\infty}u_{n}\neq0 $ ,则 $ \sum_{n=1}^{\infty}u_{n} $ 必发散.
    \item 若 $ u_{n}\geq0,u_{n}\geq u_{n+1}(n=1,2,3\ldots) $ 且 $ \lim_{n\to\infty}u_{n}=0, $ 则 $ u_{n} $ 必收敛.
    \item 若 $ \sum_{n=1}^{\infty}u_{n} $ 收敛，则 $ \sum_{n=1}^{\infty}|u_{n}| $ 必收敛.
    \item 若 $ \sum_{n=1}^{\infty}u_{n} $ 收敛于 $s$,则任意改变该级数项的位置所得到的新的级数仍收敛于 $s$.
\end{enumerate}
\choicesFour{0}{1}{2}{3}

\begin{redenv}
\analysisrule
\subsection*{逐步解析}

\begin{detailedanalysis}
\item \textbf{详细求解过程}
\begin{enumerate}[label=(\arabic*)]
    \item 正确。这是级数收敛的必要条件的逆否命题。
    \item 错误。这描述的是莱布尼茨判别法的前半部分，但缺少了交错级数的条件。对于正项级数，单调递减趋于0不一定收敛，如 $\sum 1/n$。此处语意为“数列 $u_n$ 必收敛到0”是废话因为前提给了。如果是指“级数 $\sum u_n$”收敛，则是错误的。
    \item 错误。反例：交错调和级数 $\sum (-1)^n/n$ 收敛，但取绝对值后 $\sum 1/n$ 发散。这叫条件收敛。
    \item 错误。只有绝对收敛级数才能任意交换次序保持收敛性和和不变（黎曼重排定理说明条件收敛级数重排后和可以改变）。
\end{enumerate}
只有 (1) 是正确的。

\item \textbf{最终答案}
\textbf{答案}：B
\end{detailedanalysis}
\end{redenv}

% ========== 第2题 ==========
\problemrule
\section*{【题目2】级数敛散性判定}
下列级数中收敛的级数是 \blank{1cm}。
\choicesFour{$ \sum_{n=1}^{\infty}\frac{1}{3^{n}} $}{$ \sum_{n=1}^{\infty}\frac{1}{n+1} $}{$ \sum_{n=1}^{\infty}\frac{3^{n}}{2^{n}} $}{$ \sum_{n=1}^{\infty}\frac{1}{\sqrt{n+1}} $}

\begin{redenv}
\analysisrule
\subsection*{逐步解析}

\begin{detailedanalysis}
\item \textbf{详细求解过程}
\begin{enumerate}[label=\Alph*.]
    \item 公比 $q=1/3 < 1$，收敛的几何级数。
    \item 调和级数去项，发散。
    \item 公比 $q=3/2 > 1$，发散。
    \item p-级数，$p=1/2 \le 1$，发散。
\end{enumerate}

\item \textbf{最终答案}
\textbf{答案}：A
\end{detailedanalysis}
\end{redenv}

% ========== 第3题 ==========
\problemrule
\section*{【题目3】绝对收敛判定}
级数 $ \sum_{n=1}^{\infty} (-1)^{n} \frac{1}{n^{2}} $ 是 \blank{1cm}。
\choicesFour{发散的}{条件收敛}{绝对收敛}{收敛性不能确定}

\begin{redenv}
\analysisrule
\subsection*{逐步解析}

\begin{detailedanalysis}
\item \textbf{详细求解过程}
取绝对值后为 $\sum \frac{1}{n^2}$。
这是 $p=2 > 1$ 的 p-级数，收敛。
所以原级数绝对收敛。

\item \textbf{最终答案}
\textbf{答案}：C
\end{detailedanalysis}
\end{redenv}

% ========== 第4题 ==========
\problemrule
\section*{【题目4】条件收敛判定}
下列级数中为条件收敛的级数是 \blank{1cm}。
\choicesTwo
{$ \sum_{n=1}^{\infty}(-1)^{n}\frac{n}{n+1} $}
{$ \sum_{n=1}^{\infty}(-1)^{n}\sqrt{n} $}
{$ \sum_{n=1}^{\infty}(-1)^{n}\frac{1}{n^{2}} $}
{$ \sum_{n=1}^{\infty}(-1)^{n}\frac{1}{\sqrt{n}} $}

\begin{redenv}
\analysisrule
\subsection*{逐步解析}

\begin{detailedanalysis}
\item \textbf{详细求解过程}
\begin{enumerate}[label=\Alph*.]
    \item 一般项如果不趋于0（极限为1或-1），发散。
    \item 一般项不趋于0，发散。
    \item 绝对值级数收敛，故绝对收敛。
    \item 莱布尼茨判别法：收敛。绝对值级数 $\sum 1/\sqrt{n}$ 发散。故条件收敛。
\end{enumerate}

\item \textbf{最终答案}
\textbf{答案}：D
\end{detailedanalysis}
\end{redenv}

% ========== 第5题 ==========
\problemrule
\section*{【题目5】收敛半径}
幂级数 $ \sum_{n=0}^{\infty}\frac{x^{n}}{n!} $ 的收敛半径是 \blank{1cm}。
\choicesFour{1}{0}{$ +\infty $}{$ \frac{1}{2} $}

\begin{redenv}
\analysisrule
\subsection*{逐步解析}

\begin{detailedanalysis}
\item \textbf{详细求解过程}
系数 $a_n = 1/n!$。
$\rho = \lim_{n \to \infty} |\frac{a_n}{a_{n+1}}| = \lim \frac{(n+1)!}{n!} = \lim (n+1) = +\infty$。
收敛半径 $R = +\infty$。

\item \textbf{最终答案}
\textbf{答案}：C
\end{detailedanalysis}
\end{redenv}

% ========== 第6题 ==========
\problemrule
\section*{【题目6】收敛域求解}
求幂级数 $ \sum_{n=1}^{\infty} \frac{(x-5)^{n}}{\sqrt{n}} $ 的收敛域 \blank{2cm}。

\begin{redenv}
\analysisrule
\subsection*{逐步解析}

\begin{detailedanalysis}
\item \textbf{详细求解过程}
\begin{enumerate}[label=(\arabic*)]
\item \textbf{求收敛半径}：
$a_n = 1/\sqrt{n}$。
$R = \lim |\frac{a_n}{a_{n+1}}| = 1$。
收敛区间 $(5-1, 5+1) = (4, 6)$。
\item \textbf{考察端点}：
$x=6$: $\sum \frac{1}{\sqrt{n}}$，发散。
$x=4$: $\sum \frac{(-1)^n}{\sqrt{n}}$，条件收敛。
\item \textbf{收敛域}：
$[4, 6)$。
\end{enumerate}

\item \textbf{最终答案}
\textbf{答案}：$ [4, 6) $
\end{detailedanalysis}
\end{redenv}

% ========== 第7题 ==========
\problemrule
\section*{【题目7】正项级数判别}
判别级数的收敛性，级数是 $ \sum_{n=1}^{\infty}\frac{n^{n}}{(n!)^{2}} $ \blank{2cm}。(收敛还是发散)

\begin{redenv}
\analysisrule
\subsection*{逐步解析}

\begin{detailedanalysis}
\item \textbf{详细求解过程}
比值判别法：
$$ \lim_{n \to \infty} \frac{(n+1)^{n+1}/((n+1)!)^2}{n^n/(n!)^2} = \lim \frac{(n+1)^{n+1}}{n^n} \frac{(n!)^2}{((n+1)!)^2} $$
$$ = \lim (n+1)(\frac{n+1}{n})^n \cdot \frac{1}{(n+1)^2} = \lim \frac{1}{n+1} (1+\frac{1}{n})^n = 0 \cdot e = 0 < 1 $$
所以级数收敛。

\item \textbf{最终答案}
\textbf{答案}：收敛
\end{detailedanalysis}
\end{redenv}

% ========== 第8题 ==========
\problemrule
\section*{【题目8】交错级数判别}
判别级数的收敛性，级数是 $ \sum_{n=0}^{\infty} (-1)^{n} \frac{(n+1)!}{n^{n+1}} $ \blank{2cm}。(绝对收敛，条件收敛还是发散)

\begin{redenv}
\analysisrule
\subsection*{逐步解析}

\begin{detailedanalysis}
\item \textbf{详细求解过程}
考察绝对值级数 $\sum \frac{(n+1)!}{n^{n+1}}$。
比值法：
$$ \lim \frac{(n+2)!}{(n+1)^{n+2}} \cdot \frac{n^{n+1}}{(n+1)!} = \lim (n+2) \frac{n^{n+1}}{(n+1)^{n+2}} $$
$$ = \lim \frac{n+2}{n+1} (\frac{n}{n+1})^{n+1} = 1 \cdot \frac{1}{e} = \frac{1}{e} < 1 $$
绝对值级数收敛。
所以原级数绝对收敛。

\item \textbf{最终答案}
\textbf{答案}：绝对收敛
\end{detailedanalysis}
\end{redenv}

% ========== 第9题 ==========
\problemrule
\section*{【题目9】收敛域}
求幂级数 $ \sum_{n=0}^{\infty} 10^{n}(x-1)^{n} $ 的收敛域 \blank{2cm}。

\begin{redenv}
\analysisrule
\subsection*{逐步解析}

\begin{detailedanalysis}
\item \textbf{详细求解过程}
$a_n = 10^n$。
$R = 1/10$。
中心 $x=1$。区间 $(0.9, 1.1)$。
端点处一般项不趋于0，发散。
收敛域 $(0.9, 1.1)$。

\item \textbf{最终答案}
\textbf{答案}：$ (0.9, 1.1) $ （或 $( \frac{9}{10}, \frac{11}{10} )$）
\end{detailedanalysis}
\end{redenv}

% ========== 第10题 ==========
\problemrule
\section*{【题目10】和函数}
求幂级数 $ \sum_{n=0}^{\infty} \frac{x^{n}}{2^{n-1}} $ 的和函数 \blank{2cm}。

\begin{redenv}
\analysisrule
\subsection*{逐步解析}

\begin{detailedanalysis}
\item \textbf{详细求解过程}
$$ \sum_{n=0}^{\infty} 2 \cdot (\frac{x}{2})^n = 2 \sum_{n=0}^{\infty} (\frac{x}{2})^n $$
这是几何级数，公比 $x/2$。
当 $|x/2| < 1$ 即 $|x| < 2$ 时收敛。
和 $S(x) = 2 \cdot \frac{1}{1 - x/2} = \frac{4}{2-x}$。

\item \textbf{最终答案}
\textbf{答案}：$ \frac{4}{2-x} \, (|x|<2) $
\end{detailedanalysis}
\end{redenv}

% ========== 第11题 ==========
\problemrule
\section*{【题目11】收敛域}
求幂级数 $ \sum_{n=0}^{\infty}\frac{x^{n}}{\sqrt{n+1}} $ 的收敛域。

\begin{redenv}
\analysisrule
\subsection*{逐步解析}

\begin{detailedanalysis}
\item \textbf{详细求解过程}
$R=1$。区间 $(-1, 1)$。
$x=1$: $\sum \frac{1}{\sqrt{n+1}}$ 发散。
$x=-1$: $\sum \frac{(-1)^n}{\sqrt{n+1}}$ 收敛。
收敛域 $[-1, 1)$。

\item \textbf{最终答案}
\textbf{答案}：$ [-1, 1) $
\end{detailedanalysis}
\end{redenv}

% ========== 第12题 ==========
\problemrule
\section*{【题目12】和函数与积分}
求幂级数 $ \sum_{n=0}^{\infty} \frac{x^{n+2}}{n+1} $ 的收敛域及和函数。

\begin{redenv}
\analysisrule
\subsection*{逐步解析}

\begin{detailedanalysis}
\item \textbf{详细求解过程}
\begin{enumerate}[label=(\arabic*)]
\item \textbf{收敛域}：$R=1$。
$x=1$ 发散；$x=-1$ 收敛。域 $[-1, 1)$。
\item \textbf{和函数}：
提取 $x$：$S(x) = x \sum_{n=0}^{\infty} \frac{x^{n+1}}{n+1}$。
令 $T(x) = \sum_{n=0}^{\infty} \frac{x^{n+1}}{n+1} = \sum_{k=1}^{\infty} \frac{x^k}{k} = -\ln(1-x)$。
所以 $S(x) = -x\ln(1-x)$。
\end{enumerate}

\item \textbf{最终答案}
\textbf{答案}：$ [-1, 1); S(x)=-x\ln(1-x) $
\end{detailedanalysis}
\end{redenv}

% ========== 第13题 ==========
\problemrule
\section*{【题目13】和函数与导数}
求幂级数 $ \sum_{n=1}^{\infty} n(x-1)^{n} $ 的收敛域及和函数。

\begin{redenv}
\analysisrule
\subsection*{逐步解析}

\begin{detailedanalysis}
\item \textbf{详细求解过程}
令 $t = x-1$。$\sum n t^n$。
$R=1$。域 $t \in (-1, 1)$，即 $x \in (0, 2)$。
$S(t) = t \sum n t^{n-1} = t (\sum t^n)' = t (\frac{t}{1-t})'$（注意 $n$ 从1开始，$\sum_{n=1} t^n = \frac{t}{1-t}$）。
或者：$\sum_{n=0} t^n = \frac{1}{1-t}$。求导 $\sum n t^{n-1} = \frac{1}{(1-t)^2}$。
所以 $\sum n t^n = \frac{t}{(1-t)^2}$。
代回 $x$：$\frac{x-1}{(2-x)^2}$。

\item \textbf{最终答案}
\textbf{答案}：$ (0, 2); S(x)=\frac{x-1}{(2-x)^2} $
\end{detailedanalysis}
\end{redenv}

% ========== 第14题 ==========
\problemrule
\section*{【题目14】级数求和}
求幂级数 $ \sum_{n=2}^{\infty}\frac{1}{(n^{2}-1)2^{n}} $ 的收敛域及和函数（此处是数项级数还是幂级数？题目写幂级数但无 $x$。推测是求和。题目可能是 $ \sum \frac{x^n}{(n^2-1)} $?
\emph{修正}：题目原文确实写“幂级数”，但项中无 $x$。这通常意味着求该数项级数的值。
或者 $x$ 被取为 $1/2$。
假设题目是求值。

\begin{redenv}
\analysisrule
\subsection*{逐步解析}

\begin{detailedanalysis}
\item \textbf{详细求解过程}
$$ \frac{1}{n^2-1} = \frac{1}{2}(\frac{1}{n-1} - \frac{1}{n+1}) $$
原式 $= \frac{1}{2} \sum_{n=2}^\infty (\frac{1}{n-1} - \frac{1}{n+1}) \frac{1}{2^n}$。
构造幂级数 $f(x) = \sum_{n=2}^\infty \frac{x^n}{n^2-1}$。求 $f(1/2)$。
$f(x) = \frac{1}{2} \sum (\frac{x^n}{n-1} - \frac{x^n}{n+1}) = \frac{1}{2} [x \sum \frac{x^{n-1}}{n-1} - \frac{1}{x} \sum \frac{x^{n+1}}{n+1}]$。
$\sum_{n=2} \frac{x^{n-1}}{n-1} = -\ln(1-x)$。
$\sum_{n=2} \frac{x^{n+1}}{n+1} = \sum_{k=3} \frac{x^k}{k} = -\ln(1-x) - x - x^2/2$。
$f(x) = \frac{1}{2} [-x\ln(1-x) - \frac{1}{x}(-\ln(1-x) - x - x^2/2)] = \frac{1}{2} [(\frac{1}{x}-x)\ln(1-x) + 1 + x/2]$。
代入 $x=1/2$：
$f(1/2) = \frac{1}{2} [(2 - 0.5)\ln(0.5) + 1 + 0.25] = \frac{1}{2} [1.5(-\ln 2) + 1.25] = \frac{5}{8} - \frac{3}{4}\ln 2$。

\item \textbf{最终答案}
\textbf{答案}：$ \frac{5}{8} - \frac{3}{4}\ln 2 $
\end{detailedanalysis}
\end{redenv}

% ========== 第15题 ==========
\problemrule
\section*{【题目15】幂级数求和}
求幂级数 $ \sum_{n=1}^{\infty} (-1)^{n} \frac{n(n+1)}{2^{2n}} $ 的收敛域及和函数（同样无 $x$，是求值）。

\begin{redenv}
\analysisrule
\subsection*{逐步解析}

\begin{detailedanalysis}
\item \textbf{详细求解过程}
项为 $n(n+1) x^n$ 在 $x = -1/4$ 处的值。
$S(x) = \sum n(n+1) x^n = x \sum n(n+1) x^{n-1} = x (\sum (n+1) x^n)'$。
$\sum (n+1) x^n = (\sum x^{n+1})' = (\frac{x}{1-x})'$? No.
$\sum_{n=1} x^{n+1} = \frac{x^2}{1-x}$。导数 $\frac{2x(1-x) - x^2(-1)}{(1-x)^2} = \frac{2x-x^2}{(1-x)^2}$。
再求导：$(\frac{2x-x^2}{(1-x)^2})'$。
计算量大。
替代法：$\sum x^n = \frac{1}{1-x}$。
求导两次：$\sum n(n-1) x^{n-2} = \frac{2}{(1-x)^3}$。
原式 $\sum (n^2+n) = \sum n(n-1) + \sum 2n$。
$\sum n(n+1) x^n = \sum n(n-1)x^n + \sum 2n x^n = x^2 \frac{2}{(1-x)^3} + 2x \frac{1}{(1-x)^2} = \frac{2x^2 + 2x(1-x)}{(1-x)^3} = \frac{2x}{(1-x)^3}$。
代入 $x = -1/4$：
$S = \frac{-1/2}{(5/4)^3} = \frac{-1/2}{125/64} = -\frac{32}{125}$。

\item \textbf{最终答案}
\textbf{答案}：$ -\frac{32}{125} $
\end{detailedanalysis}
\end{redenv}

\end{document}
